% concatenated from main.tex
\documentclass[11pt]{amsart}
\usepackage{epsfig,color}
\usepackage{blindtext}
\usepackage{graphicx}
\usepackage{enumitem}
\usepackage{url}
\usepackage{amssymb}
\usepackage{graphicx,import}
\usepackage{comment}
\usepackage{esint}
\usepackage{xypic}
\usepackage{amsthm}

%% page size
%\headheight=6.15pt \textheight=8in \textwidth=6.5in
%\oddsidemargin=0in \evensidemargin=0in \topmargin=0in
\usepackage[margin = 1.4in] {geometry}

\usepackage{hyperref}

%------------------------------------------------------------------------------------%
\newtheorem*{restate}{Theorem}
\newtheorem{thm}{Theorem}[section]
\newtheorem{cor}[thm]{Corollary}
\newtheorem{lem}[thm]{Lemma}
\newtheorem{prop}[thm]{Proposition}
\newtheorem{defn}[thm]{Definition}
\newtheorem{rem}[thm]{\bf{Remark}}
\newtheorem{alg}{\bf{Algorithm}}
\newtheorem*{Proof}{Proof}
\newtheorem{ex}[thm]{Example}
\numberwithin{equation}{section}
\def\pn{\par\noindent}
\def\cen{\centerline}
%------------------------------------------------------------------------------------%
% theorem environments
\setcounter{section}{0}
\theoremstyle{theorem}

\newtheorem{theorem}{Theorem}[section]
\newtheorem{conjecture}[theorem]{Conjecture}
\newtheorem{proposition}[theorem]{Proposition}
\newtheorem{lemma}[theorem]{Lemma}
\newtheorem{corollary}[theorem]{Corollary}
\newtheorem{question}[theorem]{Question}
\newtheorem{fact}[theorem]{Fact}

\theoremstyle{definition}
\newtheorem{example}[theorem]{Example}
\newtheorem{definition}[theorem]{Definition}
\newtheorem*{remark*}{Remark}
\newtheorem*{acknowledgements}{Acknowledgements}
\newtheorem{remark}[theorem]{Remark}


\numberwithin{equation}{section}


% commands
\newcommand{\cA}{\mathcal{A}}
\renewcommand{\cH}{\mathcal{H}}
\newcommand{\R}{\mathbb{R}}
\newcommand{\Q}{\mathbb{Q}}
\newcommand{\N}{\mathbb{N}}
\newcommand{\Z}{\mathbb{Z}}

\DeclareMathOperator{\Ric}{Ric}
\DeclareMathOperator{\Lip}{Lip}
\DeclareMathOperator{\diam}{diam}
\newcommand{\eps}{\varepsilon}
\DeclareMathOperator{\Image}{image}
\DeclareMathOperator{\codim}{codim}
\DeclareMathOperator{\inj}{inj}
\DeclareMathOperator{\vol}{Vol} 


\author{Elham Matinpour}
\address{Department of Mathematics, Johns Hopkins University, 3400 N. Charles Street, Baltimore, MD 21218}
\email{ematinp1@jhu.edu}


\title{Rigidity and index of free boundary minimal $Y$-cones in the unit ball}


\begin{document}


\maketitle

% ABSTRACT
\begin{abstract}
J.~C.~C.~Nitsche~\cite{nitsche1985stationary} proved that any minimal disk satisfying the free boundary condition in the unit ball $\mathbb{B}^3$ must be an equatorial flat disk. Later, Fraser and Schoen~\cite{fraser2015uniqueness} extended this rigidity theorem to higher dimensions and to ambient spaces of constant curvature. 
In this paper, we establish an analogue of the Nitsche--Fraser--Schoen theorem for singular free boundary minimal surfaces of $Y$-type in $\mathbb{B}^n$. Specifically, we prove that any conformal and minimal immersion of the standard compact flat $Y$-cone meeting $\partial \mathbb{B}^n$ orthogonally must itself be a flat $Y$-cone.
In addition, we compute the Morse index of the free boundary flat $Y$-cone and show that it equals $2(n-2)$. Furthermore, we prove that $Y$-cones are the only free boundary minimal surfaces in $\mathbb{B}^n$ with Morse index $2(n-2)$.
\end{abstract}



\section{Introduction}
A \emph{free boundary minimal surface} (FBMS) in the unit ball
$\mathbb{B}^n\subset\mathbb{R}^n$ is a minimal surface whose boundary
lies on the sphere $\mathbb{S}^{n-1}$ and meets it orthogonally.
\medskip

A classical rigidity result of Nitsche~\cite{nitsche1985stationary} states that 
any minimal disk satisfying the free boundary condition in $\mathbb{B}^3$ must 
be an equatorial flat disk. Fraser and Schoen~\cite{fraser2015uniqueness} 
extended this theorem to all dimensions, proving that a minimal disk with free 
boundary in a ball of constant curvature must be totally geodesic. They also 
showed that if the disk has parallel mean curvature vector, then it lies in a 
three–dimensional constant curvature submanifold and is totally umbilic. 
Related earlier work includes results of Souam~\cite{souam1997stability}.

Another important problem concerns the classification of free boundary minimal 
annuli. Nitsche and Fraser--Li~\cite{fraser2014compactness} conjectured that any 
embedded free boundary minimal annulus in a ball must be the critical catenoid, 
namely the portion of the catenoid that meets the boundary sphere orthogonally. We note that the uniqueness of the critical catenoid depends on whether the embeddedness condition is imposed; for related discussions on immersed free boundary minimal annuli, see Fernández--Hauswirth--Mira~\cite{fernandez2023free}.
Recent progress toward this conjecture has been made by Choe~\cite{choe2025free}, 
who developed a reflection principle for free boundary minimal surfaces, showing 
that the sphere can play a role analogous to a mirror plane.
\medskip

The purpose of this paper is to establish an analogue of the rigidity results 
of Nitsche and Fraser--Schoen in the singular setting of $Y$--type free boundary 
minimal surfaces. Such surfaces consist of three smooth faces meeting along a 
common curve with $120^\circ$ angles, forming a $Y$--configuration. Our first 
result shows that the flat $Y$--cone is the only conformal minimal immersion of 
this type in the unit ball.

\begin{theorem}\label{thm:main1}
Let $YC$ be the union of three half--unit disks joined along their diameters, 
meeting at equal angles of $120^\circ$. Suppose 
$u : YC \to \mathbb{B}^n$ is a conformal minimal immersion such that 
$u(YC)$ meets $\partial \mathbb{B}^n$ orthogonally. Then $u(YC)$ is 
congruent to the flat $Y$--cone $YC$.
\end{theorem}


We also study the stability properties of the flat free boundary $Y$--cone. 
First we compute its Morse index and nullity in the three--dimensional ball.

\begin{theorem}\label{thm:main2}
The Morse index of the free boundary flat $Y$--cone in $\mathbb{B}^3$ 
is two and its nullity is five.
\end{theorem}

The proof relies on analyzing the second variation of area for compatible
normal variations of the $Y$--surface. Since each face of the flat
$Y$--cone is planar, the Jacobi operator reduces to the Laplacian and the
stability spectrum can be computed explicitly via a Steklov eigenvalue
problem. A direct computation shows that the index form is negative definite on the 
two--dimensional space generated by constant modes.
\medskip

Finally, we extend this analysis to higher dimensions.

\begin{theorem}\label{thm:main2-n}
The Morse index of the free boundary flat $Y$--cone in $\mathbb{B}^n$ 
is $2(n-2)$ and its nullity is $5(n-2)$. Moreover, if 
$\Sigma \subset \mathbb{B}^n$ is a free boundary minimal $Y$--surface with 
Morse index $2(n-2)$, then $\Sigma$ must be the flat $Y$--cone.
\end{theorem}

The higher--dimensional result follows from the fact that the normal bundle of 
the cone has rank $n-2$. Each of the two negative modes in $\mathbb{B}^3$ 
extends independently in all normal directions, yielding index $2(n-2)$. 
Moreover, the maximality of the index forces $(n-2)$ independent ambient 
parallel vector fields to be tangent to $\Sigma$, implying that $\Sigma$ lies 
in a three--dimensional affine subspace, where the classification reduces to 
Theorem~\ref{thm:main2}.

\subsection*{Acknowledgment} I would like to thank Professor Jaigyoung Choe for suggesting the problem and Professor Jacob Bernstein for his helpful comments.


\section{Main discussion and Rigidity Result} 
Let $\Sigma_1,\Sigma_2,\Sigma_3$ be compact smooth two–dimensional surfaces with boundary.
Let $\Gamma$ be a smooth one–dimensional manifold.
For each $j=1,2,3$, suppose that the boundary of $\Sigma_j$ can be decomposed as
\[
\partial\Sigma_j=\partial'\Sigma_j\cup\partial''\Sigma_j,
\]
where $\partial'\Sigma_j$ and $\partial''\Sigma_j$ are disjoint smooth arcs (possibly sharing endpoints).

Assume there exist smooth diffeomorphisms
\[
\phi_j:\Gamma\to\partial''\Sigma_j .
\]
Identifying the boundary arcs $\partial''\Sigma_j$ via the maps $\phi_j$, we obtain a surface
\[
\Sigma=\Big(\bigcup_{j=1}^3\Sigma_j;\Gamma\Big),
\]
called a \emph{triple junction surface}. The curve $\Gamma$ is the \emph{junction curve}, along which the three faces meet. The boundary of $\Sigma$ is therefore the union of the non-identified boundary parts,
\[
\partial \Sigma = \bigcup_{j=1}^3 \partial' \Sigma_j,
\]
For clarity of notation, we use $\Gamma$ to denote the identified junction of the entire triple junction surface, and retain the notation $\partial'' \Sigma_j$ when referring to the boundary portion belonging to an individual face $\Sigma_j$ prior to identification. Thus, $\Gamma$ and $\partial'' \Sigma_j$ describe the same geometric curve viewed from different perspectives.
\medskip

Suppose that each face $\Sigma_j$ is minimally immersed in $\mathbb{R}^n$, that is,
the mean curvature vector of $\Sigma_j$ vanishes identically. Furthermore, assume that
for each $p\in\Gamma$ the outward conormal vectors of the three faces along the junction
satisfy the $Y$–balance condition
\[
\tau_1+\tau_2+\tau_3=0 .
\] 
In this case, we say that 
\[
\Sigma = \Big(\bigcup_{j=1}^3 \Sigma_j; \Gamma \Big)
\]
is a \emph{minimal $Y$-surface} in the unit ball.

\begin{definition}\label{def:minimal_immersion}
Let $\Sigma = \big(\bigcup_{j=1}^3 \Sigma_j ; \Gamma \big)$ be a triple junction surface.
A map
\[
\varphi=(\varphi_1,\varphi_2,\varphi_3):\Sigma \to \mathbb{R}^n
\]
is called a \emph{minimal immersion} if the following conditions hold:

\begin{enumerate}
\item[(i)] For each $j=1,2,3$, the map
\[
\varphi_j:\Sigma_j\to\mathbb{R}^n
\]
is a smooth minimal immersion.

\item[(ii)] Along the junction curve $\Gamma$, the three maps agree:
\[
\varphi_1|_{\partial''\Sigma_1}
=
\varphi_2|_{\partial''\Sigma_2}
=
\varphi_3|_{\partial''\Sigma_3}.
\]
Hence the images of the three boundary arcs coincide and form a common immersed curve
\[
\gamma=\varphi(\Gamma).
\]

\item[(iii)] Let $\tau_j$ denote the outward unit conormal vector to the surface
$\varphi_j(\Sigma_j)$ along $\gamma$. Then the conormals satisfy the
\emph{balance condition}
\[
\tau_1+\tau_2+\tau_3=0 \quad \text{along } \gamma .
\]
\end{enumerate}

Condition (iii) is equivalent to requiring that the three faces meet along
$\gamma$ with pairwise angles equal to $120^\circ$, forming a $Y$-configuration.
\end{definition}

Assume that both $\partial' \Sigma_j$ and $\partial'' \Sigma_j$ are nonempty. 
Suppose further that the identified boundary parts $\partial'' \Sigma_j$ form a smooth immersed one-dimensional curve in $\mathbb{B}^n$, while the non-identified boundary parts $\partial' \Sigma_j$ lie on the unit sphere $\mathbb{S}^{n-1}$. 

\begin{definition}\label{def:freeboundary}
The triple junction surface $\Sigma$ is called a \emph{free boundary minimal $Y$–surface} in $\mathbb{B}^n$ if

\begin{enumerate}
\item each face $\Sigma_j$ is minimally immersed in $\mathbb{B}^n$,
\item the boundary $\partial\Sigma=\bigcup_{j=1}^3\partial'\Sigma_j$ lies in $\mathbb{S}^{n-1}$,
\item the surfaces meet $\mathbb{S}^{n-1}$ orthogonally,
\item along the junction curve $\Gamma$ the outward conormals $\tau_1,\tau_2,\tau_3$ satisfy
\[
\tau_1+\tau_2+\tau_3=0 ,
\]
equivalently, they meet at angles $120^\circ$.
\end{enumerate}
\end{definition}


$\Sigma$ can be equipped with an abstract metric, not
necessarily induced by a Euclidean metric. Let $g_j$ be a Riemannian metric on $\Sigma_j$. 
We say that \emph{$(g_1,g_2,g_3)$ defines a metric on $\Sigma$} if the induced metrics agree along the junction:
\[
g_i|_{\partial'' \Sigma_i} = g_j|_{\partial'' \Sigma_j}, \quad \text{for all } i,j\in \{1,2,3\},
\]
and if each $\phi_j : \Gamma \rightarrow \partial'' \Sigma_j$ is an isometry with respect to this metric. 
Equivalently, we require that $g_j|_{\partial'' \Sigma_j} = \phi_j^* g_j|_{\partial'' \Sigma_j} = g_{\Gamma}$. 
In other words, $(g_1,g_2,g_3)$ defines a metric on $\Sigma$ if all $g_j$ induce the same metric on $\Gamma$. 
Furthermore, we require the metric $g$ to be compatible with the $Y$-configuration on $\Sigma$.
\medskip

Given such a metric $g$ on $\Sigma$, we define the unit outer conormal vector field $\tau_j$ to $\Sigma_j$ along $\Gamma = \partial'' \Sigma_j$. 
By convention, $\tau_j$ points outwards from $\Sigma_j$, and satisfies $g_j(\tau_j ,\tau_j)=1$ along $\Gamma$. 
Let $\eta$ be the unit tangent vector field along $\Gamma$, so that $g_j(\tau_j ,\eta )=0$. 
\medskip

Fix $\eta$ along $\Gamma$. For $p\in \Gamma$, define the geodesic curvature of $\Gamma$ in $\Sigma_j$ by
\[
\kappa_j(p)=g_j(\nabla_{\eta}^{\Sigma_j} \eta ,\tau_j ),
\]
where $\nabla^{\Sigma_j}$ denotes the \emph{covariant derivative} on $\Sigma_j$ with respect to the metric $g_j$.

\begin{definition}
We say that the metric $g$ on $\Sigma$ is \emph{compatible with the $Y$-structure} of $\Sigma$ if the sum of the geodesic curvatures of $\Gamma$ in the faces $\Sigma_j$ satisfies
\[
\sum_{j=1}^{3} \kappa_j(p) = 0 \quad \text{along } \Gamma.
\]
\end{definition}

If $\Sigma$ is a minimal $Y$-surface in $\mathbb{R}^n$ then the sum of the outer conormals $\tau_j$ to the faces $\Sigma_j$ vanishes along the set of $Y$-points $\Gamma$, which implies that the sum of the corresponding geodesic curvatures also vanishes.
\medskip


As this paper concerns minimal immersions of flat $Y$-cones, we first define such a cone explicitly. 
\begin{definition}[Flat $Y$–cone]\label{def:Ycone}
Let
\[
\hat D=\{(x,y)\in\mathbb{R}^2: x^2+y^2\le1,\ x\ge0\}
\]
be the half–disk. Let $\hat D_\theta$ denote the rotation of $\hat D$ by angle $\theta$ about the $y$–axis in $\mathbb{R}^n$.
The \emph{flat $Y$–cone} is defined by
\[
YC=\hat D\cup\hat D_{120}\cup\hat D_{-120}.
\]
\end{definition}
Set
\[
\gamma = \{ (x,y) \in \mathbb{R}^2 : x=0,\ -1 \leq y \leq 1 \}
\quad \text{and} \quad 
\sigma = \{ (x,y) \in \mathbb{R}^2 : x^2 + y^2 = 1,\ x \geq 0 \}.
\]

We equip $YC$ with the compatible metric $g=(g_1,g_2,g_3)$ induced from the ambient space. 

\begin{definition}\label{def:conformal}
A map $u=(u_1,u_2,u_3):YC\to\mathbb{R}^n$ is called \emph{conformal} if each $u_j$ is conformal with respect to the induced metric and
\[
u_1(p)=u_2(p)=u_3(p), \qquad p\in\gamma .
\]
\end{definition}

A natural question in the study of minimal $Y$-surfaces is whether there exists a minimal map $u :YC \rightarrow \mathbb{B}^n$, mapping the boundary of $YC$ to $\partial \mathbb{B}^n$, that satisfies the free boundary condition but is not conformal (in the sense of Definition~\ref{def:conformal}). 
The answer may depend on additional geometric constraints, such as whether the junction curve $\Gamma$ is simple, or whether $\Sigma$ is embedded with each face and the junction forming embedded submanifolds of $\mathbb{R}^n$. 
This question concerns the \emph{uniformization} of minimal $Y$–surfaces under free boundary conditions, where the complexity arises from the multiple conformal structures induced by the junction configuration, an issue absent in the smooth case. Below, we cover a short survey of uniformization results for minimal $Y$-surfaces:
\medskip


Mese and Yamada \cite{mese2006parameterized, mese2011local} studied the Teichmüller space of $YC$, an infinite–dimensional moduli space of conformal structures defined on the complex domain $YC$, consisting of three standard $2$-simplices glued along a common $1$-simplex $\gamma$. 
The infinite dimensionality results from the different ways the three copies of the $2$-simplices can be glued. 
They proved a local uniformization result for minimal $Y$-surfaces; see \cite[Theorem~6]{mese2011local}. 
In \cite{wang2021uniformization}, G.~Wang investigated two types of global uniformization for triple junction surfaces: \emph{strong} and \emph{weak} uniformizations. 
Strong uniformization seeks a new metric $\bar{g}=(\bar{g}_1, \bar{g}_2, \bar{g}_3)$ on a triple junction surface conformal to the original metric $g=(g_1,g_2,g_3)$, having constant Gaussian curvature on all three faces while preserving the original boundary diffeomorphisms $\phi_j : \Gamma \rightarrow \partial'' \Sigma_j$. 
In contrast, weak uniformization allows new boundary diffeomorphisms $\bar{\phi}_j$ and a conformal metric $\bar{g}$ with constant Gaussian curvature on each face and constant (possibly varying) geodesic curvature along $\Gamma$. 
Proposition~2.5 in \cite{wang2021uniformization} shows that if all faces $\Sigma_j$ are disk-type, then $\Sigma$ admits weak uniformization to a surface with constant nonnegative Gaussian curvature.
\medskip

The following theorem establishes a rigidity result for free boundary minimal $Y$–cones analogous to the classical theorem of Nitsche and its higher–dimensional extension by Fraser–Schoen.

\noindent
\textbf{Theorem~\ref{thm:main1}.}
\emph{
Let $YC$ be the union of three half-unit disks joined along their diameters, meeting at equal angles of $120^\circ$. 
Suppose $u : YC \to \mathbb{B}^n,$ is a conformal minimal immersion such that $u(YC)$ meets $\partial \mathbb{B}^n$ orthogonally. 
Then $u(YC)$ is the flat $Y$-cone $YC$.
}

\begin{proof}[Proof of Theorem~\ref{thm:main1}]
Let $YC$ be the flat $Y$-cone defined in Definition~\ref{def:Ycone}, 
that is, the union of three half-unit disks joined along their diameters and meeting at equal angles of $120^\circ$. 
Suppose that
\[
u = (u_1, u_2, u_3) : YC \to \mathbb{B}^n,
\]
is a conformal and minimal immersion such that $u(YC)$ meets $\partial \mathbb{B}^n$ orthogonally. 
We denote the image surface by 
\[
\Sigma = u(YC) = \bigcup_{j=1}^{3} \Sigma_j,
\]
where each face $\Sigma_j = u_j(\hat D)$ is a smooth minimal surface, 
and the three faces meet smoothly along the common curve 
\[
\Gamma = u(\gamma),
\]
forming a $Y$-type junction in $\mathbb{B}^n$.\\
The boundary $\partial \Sigma$ lies on the unit sphere $\mathbb{S}^{n-1}$, and the only singular points of $\Sigma$ are those along the junction curve $\Gamma$, which consists precisely of $Y$-points, where exactly three smooth sheets meet at equal $120^\circ$ angles. Note that the junction curve $\Gamma$ meets the sphere $\mathbb{S}^{n-1}$ at the endpoints of the boundary arcs $\partial'\Sigma_j$. 
Hence $\Gamma$ intersects $\mathbb{S}^{n-1}$ in exactly two points. 
If $\Gamma$ were contained entirely in the interior of $\mathbb{B}^n$, the image surface would contain a closed junction loop, which is impossible for a surface diffeomorphic to the free boundary $Y$–cone. 
\medskip

Throughout this computation, we use the Euclidean inner product on $\mathbb{R}^n$,
extended $\mathbb{C}$-bilinearly to $\mathbb{C}^n$. In local complex coordinates $z = x + i y$ on each half–disk domain $\hat{D}$, 
the components of $u_j$ satisfy
\[
(u_j)_{z\bar{z}} = 0, \qquad (u_j)_z \cdot (u_j)_z = 0, \quad j = 1,2,3,
\]
expressing that $u$ is harmonic and conformal. 
Writing the Laplacian in polar coordinates $(r,\theta)$ on $\hat{D}$, we have
\[
\Delta = \partial_{rr} + \frac{1}{r} \partial_r + \frac{1}{r^2}\partial_{\theta\theta}.
\]
Thus, the harmonicity condition $(u_j)_{z\bar{z}} = 0$ is equivalent to
\[
(u_j)_{rr} + \frac{1}{r}(u_j)_r + \frac{1}{r^2}(u_j)_{\theta\theta} = 0.
\]
Moreover, since $(u_j)_z \cdot (u_j)_z = 0$ means that the coordinate functions of $u_j$ satisfy 
$| (u_j)_x |^2 = | (u_j)_y |^2$ and $(u_j)_x \cdot (u_j)_y = 0$, 
the conformality condition is equivalent in polar coordinates to
\[
(u_j)_r \cdot (u_j)_\theta = 0,
\]
which expresses the orthogonality of the coordinate lines under $u_j$.
\medskip


The limits of the radial and angular derivatives along each sheet exist up to the junction because each $u_j$ extends smoothly to the boundary arc $\gamma$ of the half–disk. 
The three sheets do not meet smoothly across the junction; rather, they meet along $\Gamma$ with the $Y$–balance condition described in Definition~\ref{def:minimal_immersion}. 
Nevertheless, the derivatives of each individual map $u_j$ extend continuously to $\gamma$, so the polar derivatives are well defined there as limits from the interior of each sheet.
Hence all identities involving $r$– and $\theta$–derivatives remain valid at $r=0$ by continuity. \\
Along $\sigma$ the free boundary condition implies that $u_j(\sigma)\subset \mathbb{S}^{n-1}$ and the radial derivative 
${u_j}_r$ is normal to $\mathbb{S}^{n-1}$ and hence parallel to $u_j$.
Thus
\[
{u_j}_r = f\,u_j
\]
for some scalar function $f$ on $\sigma$. Differentiating tangentially along
$\sigma$ gives
\[
{u_j}_{r\theta} = f_\theta u_j + f\,{u_j}_\theta .
\]
Both terms are tangent to $\Sigma_j$ along $\sigma$, hence the normal
component vanishes and therefore
\begin{equation}
\label{eq:perp}
{u_j}_{r\theta}^{\perp}=0
\qquad \text{along }\sigma .
\end{equation}
Since $u_j$ is harmonic and conformal, it follows from the classical identity
for conformal minimal immersions that\\
\emph{Claim.} The quadratic differential
\[
({u_j}_{zz}^{\perp})^2
\]
is holomorphic on $\hat D$.\\
\emph{Proof of Claim:} Since $u_j$ is harmonic, we have $(u_j)_{z\bar z}=0$, and therefore
\[
({u_j}_{zz} \cdot {u_j}_{zz})_{\bar{z}} = 0.
\]
Moreover, the conformality condition gives
\[
{u_j}_{zz} \cdot {u_j}_z = \tfrac{1}{2}({u_j}_z \cdot {u_j}_z)_z = 0.
\]
Because each face $\Sigma_j$ is smoothly immersed, 
the normal bundle of $\Sigma_j$ is well-defined and smooth. 
Hence, we may decompose
\[
{u_j}_{zz}
= {u_j}_{zz}^{\perp}
+ \frac{{u_j}_{zz} \cdot {u_j}_z}{|{u_j}_z|^2}\,{u_j}_z
+ \frac{{u_j}_{zz} \cdot {u_j}_{\bar{z}}}{|{u_j}_{\bar{z}}|^2}\,{u_j}_{\bar{z}},
\]
where ${u_j}_{zz}^{\perp}$ denotes the component of ${u_j}_{zz}$ orthogonal to the tangent plane of 
$\Sigma_j = u_j(\hat{D})$.
Since the last two terms on the right hand side are tangential parts we have 
 $$
 {u_j}_{zz}^{\perp} \cdot \frac{{u_j}_{zz}\cdot {u_j}_{\bar{z}}}{\vert {u_j}_{\bar{z}}\vert^2} {u_j}_{\bar{z}} =0.
 $$
 On the other hand, ${u_j}_{\bar{z}}\cdot {u_j}_{\bar{z}}=0$, and thus we find that, 
 $$
 {u_j}_{zz}^{\perp} \cdot {u_j}_{zz}^{\perp} = {u_j}_{zz} \cdot {u_j}_{zz},
 $$
 which implies $({u_j}_{zz}^{\perp})^2={u_j}_{zz}^{\perp} \cdot {u_j}_{zz}^{\perp}$ is holomorphic. \hfill $\square$
 
 In the polar coordinates, on $\hat{D}$:
 $$
 {u_j}_{zz} =\frac{1}{4} e^{-2i\theta}\Big[  {u_j}_{rr}-\frac{{u_j}_r}{r} - \frac{1}{r^2} {u_j}_{\theta \theta} - i(\frac{2{u_j}_{r\theta}}{r}-\frac{2{u_j}_{\theta}}{r^2}) \Big],
 $$
 and,
 $$
{u_j}_{zz}^{\perp} =\frac{1}{4} e^{-2i\theta}\Big[  {u_j}_{rr}^{\perp} - \frac{1}{r^2} {u_j}_{\theta \theta}^{\perp} - \frac{2i}{r} {u_j}_{r\theta}^{\perp} \Big].
$$
Note that, on $\gamma$ 
$$
u_1(p)=u_2(p)=u_3(p),
$$
and,
\begin{equation} \label{eqns: p.derivs(r,rr)}
    \begin{split}
       &  {u_1}_{r}(p) = {u_2}_{r}(p) = {u_3}_{r}(p),  \\
       & {u_1}_{rr}(p) = {u_2}_{rr}(p) = {u_3}_{rr}(p).
    \end{split}
\end{equation}
Using equations (\ref{eqns: p.derivs(r,rr)}) in $\Delta u_j ={u_j}_{rr}+\frac{1}{r}{u_j}_r+\frac{1}{r^2}{u_j}_{\theta \theta}=0$ we find that
\begin{equation} \label{eqn: p.deriv(theta,theta)}
  {u_1}_{\theta \theta}(p) ={u_2}_{\theta \theta}(p) ={u_3}_{\theta \theta} (p),  
\end{equation}
along $\gamma$.
Moreover, by the $Y$-configuration along $\gamma$ we also have
\begin{equation}\label{eqn:p.deriv(theta)}
 {u_1}_{\theta}(p) +  {u_2}_{\theta}(p) + {u_3}_{\theta}(p) =0.
\end{equation}
 Set $Q_j(z)={u_j}_{zz}(z)$, and note that ${Q_j^{\perp}}^2(z)=Q_j^2(z)$. Define
$$
h(z)=\underset{j=1}{\overset{3}{\sum}}Q_j^2(z)=\underset{j=1}{\overset{3}{\sum}}{Q_j^{\perp}}^2(z),
$$
on $\hat{D}$. Compute
\begin{equation}
\nonumber
\begin{split}
h(z)&=\underset{j=1}{\overset{3}{\sum}}Q_j^2(z)=\underset{j=1}{\overset{3}{\sum}} \Big ( \frac{1}{16}e^{-4i\theta}   \big [{u_j}_{rr}-\frac{{u_j}_r}{r} -\frac{1}{r^2}{u_j}_{\theta \theta}-i(\frac{2{u_j}_{r\theta}}{r}-\frac{{2u_j}_{\theta}}{r^2}) \big ]^2 \Big )\\
&= \frac{e^{-4i\theta}}{16} \underset{j=1}{\overset{3}{\sum}} \Big( ({u_j}_{rr}-\frac{{u_j}_r}{r} -\frac{1}{r^2}{u_j}_{\theta \theta} )^2 -(\frac{2{u_j}_{r\theta}}{r}-\frac{2{u_j}_{\theta}}{r^2})^2\\
&\ -2i (\frac{2{u_j}_{r\theta}}{r}-\frac{2{u_j}_{\theta}}{r^2}) ({u_j}_{rr}-\frac{{u_j}_r}{r}-\frac{1}{r^2}{u_j}_{\theta \theta} ) \Big),
\end{split}
\end{equation}
and similarly,
\begin{equation}
    \nonumber
\begin{split}
h(z)& =\underset{j=1}{\overset{3}{\sum}}{Q_j^{\perp}}^2(z)=  \frac{1}{16} e^{-4i\theta} \underset{j=1}{\overset{3}{\sum}}\Big[  {u_j}_{rr}^{\perp} - \frac{1}{r^2} {u_j}_{\theta \theta}^{\perp} - \frac{2i}{r} {u_j}_{r\theta}^{\perp} \Big]^2 \\ 
&=\frac{e^{-4i\theta}}{16} \underset{j=1}{\overset{3}{\sum}} \Big( ({u_j}_{rr}^{\perp}-\frac{1}{r^2}{u_j}_{\theta \theta}^{\perp} )^2 -\frac{4}{r^2}{{u_j}_{r\theta}^{\perp}}^2-\frac{4i}{r} {u_j}_{r\theta}^{\perp} ({u_j}_{rr}^{\perp}-\frac{1}{r^2}{u_j}_{\theta \theta}^{\perp} ) \Big).
\end{split}
\end{equation}
Set $H(z)=z^4h(z)$. The imaginary part of $H(z)$ is
\begin{equation}
    \nonumber
    \begin{split}
        \operatorname{Im} H(z)&=\underset{j=1}{\overset{3}{\sum}}-\frac{ir^4}{8} (\frac{2{u_j}_{r\theta}}{r}-\frac{2{u_j}_{\theta}}{r^2}) ( {u_j}_{rr}-\frac{{u_j}_r}{r} -\frac{1}{r^2}{u_j}_{\theta \theta} )\\
        &=\underset{j=1}{\overset{3}{\sum}}-\frac{ir^3}{4} {u_j}_{r\theta}^{\perp} ({u_j}_{rr}^{\perp}-\frac{1}{r^2}{u_j}_{\theta \theta}^{\perp} ).
 \end{split}
\end{equation}
Since ${u_j}_{r\theta}^{\perp}=0$ on $\sigma$, see (\ref{eq:perp}), $H(z)$ is real on $\sigma$. It follows from equations (\ref{eqns: p.derivs(r,rr)}) and (\ref{eqn: p.deriv(theta,theta)}) that the terms ${u_j}_{rr}-\frac{{u_j}_r}{r} -\frac{1}{r^2}{u_j}_{\theta \theta}$ in $\operatorname{Im} H(z)$ are independent of $j$ along $\gamma$. By the $Y$-configuration, see (\ref{eqn:p.deriv(theta)}), the terms $(\frac{2{u_j}_{r\theta}}{r}-\frac{2{u_j}_{\theta}}{r^2})$ sum up to zero along $\gamma$, and thus
 $\operatorname{Im} H(z)$ is equal to zero along $\gamma$.\\
Therefore, the map $H(z)$ is holomorphic on $\hat{D}$ and is real on $\gamma \cup \sigma$. Hence, it must be constant on $\hat{D}$, and since it vanishes at zero, it must vanish on $\hat{D}$. By checking on the real and imaginary parts of $H(z)$, this implies that ${u_j}_{r\theta}^{\perp}=0$ and  ${u_j}_{rr}^{\perp}-{u_j}_{\theta \theta}^{\perp}=0$ on $\sigma$.\\
By the minimality, this implies that ${u_j}_{rr}^{\perp}={u_j}_{\theta \theta}^{\perp}=0$ on $\sigma$. Therefore, the second fundamental form of $\Sigma_j$ is zero on $\partial' \Sigma_j$. It follows that the second fundamental form of $\partial' \Sigma_j$ in $\mathbb{S}^{n-1}$ is zero. Hence, $\partial' \Sigma_j$ is a piece of a great circle, and thus $\partial \Sigma$ is composed of pieces of great circles.\\

Suppose that the boundary circles corresponding to $u_j$ lie in the planes $P_j$, and let $n_j$ be the unit normal vector to the plane $P_j$. Then $n_j\cdot u_j$ is a harmonic function that is equal to zero along $\sigma$. Moreover, by the free boundary condition, $\nabla_{\tau_j}(n_j\cdot u_j)=0$ along $\sigma$, where $\tau_j$ is the outward conormal vector to $\Sigma_j$ along $\sigma$. Then by the Calderon unique continuation of elliptic equations, $n_j\cdot u_j\equiv 0$ on $\hat{D}$, see \cite{gilbarg1977elliptic} for a reference. Hence, ${u_j}_{zz}={u_j}_{zz}^{\perp}=0$ for $j=1,2,3$.\\ 
Then, by ${u_j}_{zz}=0$ we find that ${u_j}_{rr}-\frac{{u_j}_r}{r}-\frac{1}{r^2}{u_j}_{\theta \theta}=0$, and then since $\Delta u_j ={u_j}_{rr}+\frac{1}{r}{u_j}_r+\frac{1}{r^2}{u_j}_{\theta \theta}=0$ we find that ${u_j}_{rr}=0$, and by ${u_j}_{zz}^{\perp}=0$ we find that ${u_j}_{r\theta}^{\perp}=0$ and ${u_j}_{rr}^{\perp}-\frac{1}{r^2}{u_j}_{\theta \theta}^{\perp}=0$, on $\hat{D}$. Thus
\[
(u_j)_{r\theta}^{\perp}
=
(u_j)_{rr}^{\perp}
=
(u_j)_{\theta\theta}^{\perp}
=0
\quad\text{on }\hat D .
\]
Hence the second fundamental form of each face $\Sigma_j$ vanishes, and therefore $u(YC)$ is flat. 
In particular the image of $\gamma$ is a straight line segment in $\mathbb{B}^n$. 
By the free boundary condition this segment meets $\mathbb{S}^{n-1}$ orthogonally, hence it is a diameter. 
Reflecting across this diameter gives a free boundary minimal disk, which by the theorem of Fraser–Schoen \cite{fraser2015uniqueness} must be an equatorial disk. 
Therefore each $\Sigma_j$ is flat and $u(YC)$ coincides with the flat $Y$–cone.  
\end{proof}


\section{Morse Index and Nullity of the Flat Y-Cone}
\begin{theorem}\label{thm:main2}
    The Morse index of the free boundary flat $Y$-cone in the unit ball $\mathbb{B}^3$ is two, and its nullity is five. Moreover, if $\Sigma$ is a free boundary minimal $Y$-surface in the unit ball with Morse index two, then $\Sigma$ is a $Y$-cone. 
\end{theorem}
\begin{proof}
Let  $\Sigma = (\bigcup_{j=1}^3 \Sigma_j; \Gamma)$ be the free boundary flat $Y$-cone $YC$ in $\mathbb{B}^3$, as described above \ref{def:Ycone}. Set $\nu=(\nu_{\Sigma_j})_{j=1}^3$ be a choice of unit normal vector field on $\Sigma$ that is compatible with the $Y$-configuration on $\Sigma$, where $\nu_{\Sigma_j}$ is a unit normal vector field on $\Sigma_j$. Let $\tau_j$ denote the outward conormal vector along 
$\partial \Sigma_j = \partial' \Sigma_j \cup \partial'' \Sigma_j$, so that
\[
\tau_j \perp \mathbb{S}^2 \text{ along } \partial' \Sigma_j, 
\qquad 
\tau_1 + \tau_2 + \tau_3 = 0 \text{ along } \partial'' \Sigma_j.
\]
The first condition corresponds to the free boundary condition, and the second expresses the $Y$-balance condition along the junction curve $\Gamma$.
\medskip

Define 
$$W_{com}^{1,2}(\Sigma) := \{ f=(f_1,f_2,f_3)\ :\ f_j  \in W^{1,2}(\Sigma_j)\ \& \ f_1+f_2+f_3=0\ \text{along}\  \Gamma \}.$$
For a single smooth face $\Sigma_j$, the natural functional space is the Sobolev space $W^{1,2}(\Sigma_j)$. 
The space $W^{1,2}_{com}(\Sigma)$ defined above consists of triples 
$f=(f_1,f_2,f_3)$ satisfying $f_j\in W^{1,2}(\Sigma_j)$ and the $Y$-compatibility condition 
$f_1+f_2+f_3=0$ along the common junction curve $\Gamma = \partial''\Sigma_j$.\\
For a normal variation $V = f\nu$ with $f \in W_{com}^{1,2}(\Sigma)$, the second variation of area is given by
\begin{equation} \label{eqn: 2nd variation fromula}
\begin{split}
   Q(f , f) = \frac{d^2}{dt^2} Area (\Sigma (t))\vert_{t=0} &=  \underset{j=1}{\overset{3}{\sum}} \Big ( \int_{\Sigma_j} (\vert \nabla_{\Sigma_j} f_j \vert^2 - \vert A_{\Sigma_j} \vert^2 f_j^2)  \\
   & +  \int_{\partial' \Sigma_j  }    (\bold{H}_{\partial' \Sigma_j} \cdot \tau_j) (f_j \vert_{\partial' \Sigma_j} )^2   - \int_{\partial'' \Sigma_j  }    (\bold{H}_{\partial'' \Sigma_j} \cdot \tau_j) (f_j \vert_{\partial'' \Sigma_j} )^2 \Big) 
   \end{split}
\end{equation}
where $\nabla_{\Sigma_j}$ denotes the gradient (Levi–Civita connection) on the surface $\Sigma_j$, $\vert A_{\Sigma_j} \vert^2$ is the squared norm of the second fundamental form of $\Sigma_j$, and $\bold{H}_{\partial' \Sigma_j}$ is the mean curvature vector of the curve $\partial' \Sigma_j$. \\
\label{def: Morse Index}The Morse index of $\Sigma$ is the maximal dimension of a subspace of $W_{com}^{1,2}(\Sigma)$ on which $Q$ is negative definite, and the nullity is the dimension of $\ker Q$ 
(see \cite{fraser2016sharp}, \cite{tran2016index}, \cite{wang2022curvature}, \cite{matinpour2024morse}).\\
On each face $\Sigma_j$, the Jacobi operator is $J_j = \Delta_{\Sigma_j} + |A_{\Sigma_j}|^2$. 
Since $\Sigma_j$ is flat, $A_{\Sigma_j} \equiv 0$, and hence $J_j = \Delta_{\Sigma_j}$. 
Moreover, $\mathbf{H}_{\partial'' \Sigma_j} = 0$ along $\Gamma$, so
\[
Q(f,f)
= \sum_{j=1}^3 \left(
\int_{\Sigma_j} - f_j J_j f_j
+ \int_{\partial \Sigma_j} f_j \frac{\partial f_j}{\partial \tau_j}
+ \int_{\partial' \Sigma_j} (\mathbf{H}_{\partial'\Sigma_j}\!\cdot\!\tau_j) f_j^2
\right).
\]
To evaluate the remaining terms in index form, we notice that the Jacobi operator has no negative eigenvalues on the faces $\Sigma_j$ (because $\Sigma_j$ are flat), so we are left to check the kernel of the Jacobi operator on $\Sigma_j$ and then to compute the corresponding Steklov eigenvalues, $\sigma_j$, along $\partial' \Sigma_j$, refer to \cite{matinpour2024morse} for more details.
\medskip

To analyze the index form, we consider the associated Jacobi--Steklov
eigenvalue problem on the $Y$-surface $\Sigma$.

Let $f=(f_1,f_2,f_3)\in W^{1,2}_{com}(\Sigma)$ with
$f_j\in W^{1,2}(\Sigma_j)$ and satisfying the compatibility condition
\[
f_1+f_2+f_3=0 \qquad \text{along }\Gamma .
\]
We say that $f$ is a $J$--Steklov eigenfunction with eigenvalue
$\sigma\in\mathbb{R}$ if
\begin{equation}\label{eq:steklov_correct}
\begin{cases}
J_j f_j =0 & \text{in } \Sigma_j, \\[4pt]
\dfrac{\partial f_j}{\partial \tau_j}=\sigma f_j
& \text{on } \partial'\Sigma_j , \\[8pt]
\dfrac{\partial f_j}{\partial \tau_j}
=
\dfrac{\partial f_i}{\partial \tau_i}
& \text{on } \Gamma \quad \text{for all } i, j.
\end{cases}
\end{equation}
The last condition arises naturally from integration by parts on the
three sheets of $\Sigma$. Indeed, for admissible test functions
$v=(v_1,v_2,v_3)\in W^{1,2}_{com}(\Sigma)$ we have
$v_1+v_2+v_3=0$ along $\Gamma$. Consequently,
\[
\int_{\Gamma}\sum_{i=1}^3 v_i
\frac{\partial f_i}{\partial \tau_i}\,ds=0
\quad\text{for all admissible } v
\]
if and only if
\[
\frac{\partial f_1}{\partial \tau_1}
=
\frac{\partial f_2}{\partial \tau_2}
=
\frac{\partial f_3}{\partial \tau_3}
\quad \text{on }\Gamma .
\]

In the case of the flat $Y$--cone each face $\Sigma_j$ is planar,
so $A_{\Sigma_j}\equiv0$ and hence the Jacobi operator reduces to
$J_j=\Delta_{\Sigma_j}$. The problem \eqref{eq:steklov_correct}
therefore defines a self--adjoint Steklov problem with discrete
spectrum
\[
0=\sigma_0<\sigma_1\le\sigma_2\le\cdots\nearrow\infty .
\]

Since the three faces are congruent, the Jacobi operators on
$\Sigma_j$ coincide; we therefore denote them simply by $J$.
\medskip

In particular, if $g$ is a Steklov eigenfunction associated to $\sigma$, 
then for any vector $\bold{c}=(c_1,c_2,c_3)\in\mathbb{R}^3$ satisfying $c_1+c_2+c_3=0$, 
there exists a unique compatible function $f_{\bold{c}}\in W^{1,2}_{com}(\Sigma)$ such that
\[
J f_{\bold{c}} = 0, \qquad 
f_{\bold{c}}|_{\partial'\Sigma_j} = c_j g, \qquad
\frac{\partial f_{\bold{c}}}{\partial \tau_j}\big|_{\partial'\Sigma_j} = \sigma c_j g,
\qquad
\frac{\partial f_{\bold{c}}}{\partial \tau_j}\big|_{\partial''\Sigma_j}=\frac{\partial f_{\bold{c}}}{\partial \tau_i}\big|_{\partial''\Sigma_i}.
\]
The existence of such Steklov eigenfunctions is standard if the kernel of the Jacobi operator for the Steklov problem is trivial. Here, $f$ is in the kernel of $J$ for the Steklov problem if $f\vert_{\partial' \Sigma_j =0}$ and $Jf=0$. We will find that the kernel of the Jacobi operator for the Steklov problem is trivial on the flat faces of $\Sigma$. Note that if $g$ is a Steklov eigenfunction for $J$, then the corresponding eigenvalue $\sigma$ is the same on all three identical faces. Let $g\in C^1(\partial' \Sigma_j)$ be a Steklov eigenfunction of $J$, $\bold{c}=(c_1,c_2,c_3)\in \mathbb{R}^3$ with $c_1+c_2+c_3=0$, and $f_{\bold{c}}\in W^{1,2}_{com}(\Sigma_j)$ the unique function satisfying
\[
Jf_{\bold{c}}=0, \quad 
f_{\bold{c}}\vert_{\partial' \Sigma_j}=c_j g, \quad 
\partial_{\tau_j}f_{\bold{c}}\vert_{\partial' \Sigma_j} = \sigma f_{\bold{c}}\vert_{\partial' \Sigma_j} = \sigma c_j g.
\]
Recall that $\partial \Sigma_j = \partial'' \Sigma_j \cup \partial' \Sigma_j$, where $\partial'' \Sigma_j = \Gamma$. 
For clarity, when we write $f_{\bold{c}}\vert_{\Gamma}^j$, we mean the restriction of $f_{\bold{c}}$ to the portion of the junction curve $\Gamma$ as viewed from the face $\Sigma_j$, 
that is, $f_{\bold{c}}\vert_{\Gamma}^j := f_{\bold{c}}\vert_{\partial'' \Sigma_j}$.\\
Then we compute the index form for such $f_{\bold{c}}$ as follows:
\begin{equation}
\nonumber
\begin{split}
Q(f_{\bold{c}} , f_{\bold{c}}) 
 & = \underset{j=1}{\overset{3}{\sum}} 
 \Bigg( \int_{\partial \Sigma_j} (f_j\vert_{\partial \Sigma_j} \frac{\partial f_j}{\partial \tau_j})  
 + \int_{\partial' \Sigma_j  }    \bold{H}_{\partial' \Sigma_j} \cdot \tau_j (f_j \vert_{\partial' \Sigma_j} )^2 \Bigg) \\
 &= \underset{j=1}{\overset{3}{\sum}} 
 \Bigg( \int_{\Gamma } f_{\bold{c}}\vert_{\Gamma}^j 
\frac{\partial f_{\bold{c}}}{\partial\tau_j}
 + \int_{\partial' \Sigma_j  } (\sigma +\bold{H}_{\partial' \Sigma_j} \cdot \tau_j) (f_{\bold{c}} \vert_{\partial' \Sigma_j} )^2  \Bigg).
\end{split}
\end{equation}
By [Proposition 4.2., \cite{matinpour2024morse}], the kernel of the Jacobi operator on $\Sigma_j$ is spanned by the linear combinations of the functions $w_0, w_1,\cdots $,
  \begin{equation}
\begin{cases}
\nonumber
& w_0 = 1,  \\
& w_n(r, \theta ) = a_n r^n \cos n\theta  + b_n r^n \sin n\theta,\ \  n\geq 1,\ a_n,b_n\in \mathbb{R},\  r\in [0,1],\ \theta \in [0,\pi]
\end{cases}
\end{equation}
where the restriction to $\partial' \Sigma_j$, $w_k\vert_{r=1}$, is the $J$-Steklov eigenfunction associated with the eigenvalue $\sigma_k$:
 \begin{equation}
\begin{cases}
\nonumber
& w_0\vert_{r=1} = 1, \hspace{6cm}  \sigma_0 = 0 \\
& w_n(r, \theta )\vert_{r=1} = a_n \cos n\theta  + b_n  \sin n\theta    \hspace{1.8cm}  \sigma_n  = n,\ n\geq 1
\end{cases}
\end{equation} 
Since $\Sigma_j$ is a half unit disk, the junction curve
$\Gamma=\partial''\Sigma_j$ corresponds to the diameter of the disk,
which is given in polar coordinates by $\theta=0$ and $\theta=\pi$. Since $\Gamma$ corresponds to the diameter of the half disk,
it is given by $\theta=0$ and $\theta=\pi$.
For $n\ge1$ we have
\[
w_n(r,0)=a_n r^n,\qquad
w_n(r,\pi)=(-1)^n a_n r^n ,
\]
while the tangential derivative along $\Gamma$ vanishes.
Consequently the boundary term
$\int_{\Gamma} f_{\bold c}\partial_{\tau_j}f_{\bold c}$
vanishes for all modes $w_n$. Note that $\bold{H}_{\partial' \Sigma_j} \cdot \tau_j = -1$, so we can simplify
$$
 Q(f_{\bold{c}} , f_{\bold{c}}) 
 = \underset{j=1}{\overset{3}{\sum}} \Big (  
 \int_{\partial' \Sigma_j  }   (\sigma - 1) (f_{\bold{c}} \vert_{\partial' \Sigma_j} )^2  \Big ) 
$$
The index form is negative definite on the subspace generated by
the constant mode $w_0$.  Because admissible functions satisfy the
compatibility condition $f_1+f_2+f_3=0$ along $\Gamma$, the space of
constant triples $(c_1,c_2,c_3)$ with $c_1+c_2+c_3=0$ is
two--dimensional.  Hence the Morse index equals two.

For the first nontrivial modes $w_1$, each face contributes a
two--dimensional space spanned by $\cos\theta$ and $\sin\theta$.
After imposing the compatibility condition on the three faces,
one obtains a five--dimensional space on which the index form
vanishes.  Thus the nullity equals five. These Jacobi fields correspond geometrically to infinitesimal tilting
of the planar faces that preserve the $Y$--balance along the junction.
%Hence the nullity of the flat free boundary $Y$--cone equals five.

For all higher modes $w_n$ with $n\ge2$, the factor $(\sigma-1)$
is positive, and therefore the index form is positive definite
on the corresponding subspaces.\\
%The five--dimensional null space also admits a natural geometric
%interpretation.  Recall that on each face $\Sigma_i$ the kernel of the
%Jacobi operator is spanned by the functions %$w_0,w_1,\ldots$, where
%$w_1(r,\theta)=a r\cos\theta + b r\sin\theta$.  The functions $w_1$
%correspond to the first nontrivial Steklov eigenfunctions on the
%boundary and form a two--dimensional space on each face.

%For the $Y$--surface $\Sigma=\bigcup_{i=1}^3\Sigma_i$, a normal
%variation is given by a triple %$(\phi_1,\phi_2,\phi_3)$ with
%$\phi_i\in C^\infty(\Sigma_i)$ satisfying the compatibility condition
%$\phi_1+\phi_2+\phi_3=0$ along the junction curve $\Gamma$.
%Consequently, the $w_1$--modes give rise to a six--dimensional space of
%Jacobi fields before imposing compatibility.  Along the junction
%$\Gamma$ we have $\theta=0$ or $\pi$, so the traces of $w_1$ reduce to
%multiples of $\cos\theta$, yielding a one--dimensional space on the
%boundary.  The compatibility condition therefore imposes a single
%linear relation among the six parameters.  As a result, the space of
%compatible $w_1$--modes has dimension $6-1=5$.

\medskip

Now, suppose that 
\[
\Sigma = \Big(\bigcup_{j=1}^{3}\Sigma_j;\,\Gamma \Big)
\]
is a free boundary minimal $Y$-surface in the unit ball $\mathbb{B}^3$ with Morse index two. 
Recall that the Morse index of $\Sigma$ is the maximal dimension of a subspace of normal variations along which the second variation of area is negative.\\
Let $T\mathbb{R}^3$ denote the tangent bundle of the ambient space
$\mathbb{R}^3$.
 Each fiber $T\mathbb{R}^3_p$ is a three--dimensional vector space spanned by the standard tangent vector fields $E_1,E_2,E_3$. 
Let $N\Sigma$ be the normal bundle of $\Sigma$, and $\Gamma(N\Sigma)$ the space of its smooth sections. 
For each ambient tangent field $E_j$, write $NE_j\in \Gamma(N\Sigma)$ for its normal projection along $\Sigma$.\\
Along the junction curve $\Gamma$, the normal bundle $N\Sigma$ is understood
in the natural piecewise sense: each face $\Sigma_i$ carries its own unit normal field $N_i$, smooth up to $\Gamma$, and the collection $\{N_1,N_2,N_3\}$ satisfies the $Y$-balance condition
$N_1+N_2+N_3=0$ on $\Gamma$.  Thus, a section of $N\Sigma$ is given by
a triple of normal fields $(\phi_1 N_1,\phi_2 N_2,\phi_3 N_3)$
with $\phi_i\in C^\infty(\Sigma_i)$ satisfying the compatibility condition
$\phi_1+\phi_2+\phi_3=0$ along $\Gamma$.
For an ambient constant vector field $E$, its normal part along $\Sigma$ is defined
facewise by $n_i=\langle E,N_i\rangle$ on $\Sigma_i$.
Since each $N_i$ extends continuously to $\Gamma$, the functions $n_i$ do as well,
and the triple $(n_1,n_2,n_3)$ satisfies $n_1+n_2+n_3=0$ along $\Gamma$.
This defines the normal projection $NE$ continuously across the junction.
\medskip

On each face $\Sigma_i$, the normal component $n_i := \langle E, N_i\rangle$ of an ambient tangent field $E$ satisfies
\[
J_i n_i = \Delta_{\Sigma_i}n_i + |A_{\Sigma_i}|^2 n_i = 0,
\]
since $E$ is tangent and $\Sigma_i$ is minimal; see \cite{colding2011course}.
Hence the volume terms in the index form vanish:
\[
\int_{\Sigma_i} -\,n_i J_i n_i = 0.
\]
Substituting into the index form~\eqref{eqn: 2nd variation fromula}, we obtain
\[
Q(n,n)
= \sum_{i=1}^3 \Big( \int_{\partial\Sigma_i} n_i\,\frac{\partial n_i}{\partial \tau_i}
+ \int_{\partial'\Sigma_i} (\boldsymbol{H}_{\partial'\Sigma_i}\!\cdot\!\tau_i)\,n_i^2
- \int_{\Gamma} (\boldsymbol{H}_{\Gamma}\!\cdot\!\tau_i)\,n_i^2 \Big),
\]
where $\partial\Sigma_i=\partial'\Sigma_i\cup\Gamma$, the curves $\partial'\Sigma_i$ lie on $\mathbb{S}^2$, and $\tau_i$ denotes the outward conormal to $\partial\Sigma_i$ in $\Sigma_i$.

We analyze these boundary terms one by one.\\
Along the common curve $\Gamma$, the conormals $\tau_i$ and the geodesic curvatures $\kappa_i := \boldsymbol{H}_{\Gamma}\!\cdot\!\tau_i$ of $\Gamma$ in each face $\Sigma_i$ satisfy the compatibility conditions
\[
\tau_1 + \tau_2 + \tau_3 = 0,
\qquad
\kappa_1 + \kappa_2 + \kappa_3 = 0,
\qquad
N_1 + N_2 + N_3 = 0.
\]
The last identity implies $n_1 + n_2 + n_3 = 0$ along $\Gamma$.
Observe that
\[
\int_{\partial\Sigma_i} n_i\frac{\partial n_i}{\partial \tau_i}
= \int_{\partial'\Sigma_i} n_i\frac{\partial n_i}{\partial \tau_i}
+ \int_{\Gamma} n_i\frac{\partial n_i}{\partial \tau_i}.
\]
 Summing over \(i\) and using these identities, one performs a local computation (see Remark~\ref{rmk:local comp. on Gamma}) to show the pointwise identity
\[
\sum_{i=1}^3 \big(n_i \partial_{\tau_i} n_i - \kappa_i n_i^2\big)
= \tfrac{1}{2}\,\partial_\eta \!\Big(\sum_{i=1}^3 \tau_i[n_i^2]\Big),
\]
where $\eta$ denotes the unit tangent to $\Gamma$. 
Integrating this expression along $\Gamma$ yields zero on the right-hand side (a total derivative in $\eta$), showing that the conormal-flux and geodesic-curvature terms cancel exactly:
\[
\sum_{i=1}^3 \int_\Gamma \big(n_i \partial_{\tau_i} n_i - \kappa_i n_i^2\big) = 0.
\]
Consequently, all boundary contributions from $\Gamma$ vanish after summation over $i$.\\
Combining the previous steps gives
\[
Q(n,n)
= \sum_{i=1}^3 \int_{\partial'\Sigma_i} 
(\boldsymbol{H}_{\partial'\Sigma_i}\!\cdot\!\tau_i)\,n_i^2.
\]
Because $\boldsymbol{H}_{\partial'\Sigma_i}\!\cdot\!\tau_i=-1$ on $\mathbb{S}^2$, we have $Q(n,n)<0$ unless $n_i\equiv 0$ for every $i$. 
Hence, any ambient parallel field $E$ whose normal projection $NE$ does not vanish identically provides a direction that decreases the area of~$\Sigma$.\\
Suppose the Morse index of $\Sigma$ equals two. 
Since the three parallel fields $E_1,E_2,E_3$ give rise to at most three independent negative directions $NE_j$, there must exist (after possibly relabeling coordinates) some $j$ for which $NE_j$ vanishes identically on~$\Sigma$. 
This means that $E_j$ is everywhere tangent to $\Sigma$, and hence tangent to each face $\Sigma_i$ and to the junction curve~$\Gamma$.\\
Since $E_j$ is a parallel vector field tangent to $\Gamma$,
the curve $\Gamma$ must lie in an affine line parallel to $E_j$.
Because $\Gamma$ is contained in the unit ball and is symmetric
with respect to the origin, it follows that $\Gamma$ is a straight
line segment through the origin. Hence, each face $\Sigma_i$ is a planar disk that meets the unit sphere orthogonally. 
Consequently, $\Sigma$ must coincide with the flat $Y$-cone.
\end{proof}

This computation reveals an interesting pattern in the Morse index stratification of free boundary minimal surfaces in $\mathbb{B}^3$. While the smooth theory yields index-one (equatorial disk) and index-four (critical catenoid) examples \cite{devyver2019index, tran2016index}, our work shows that singular $Y$-surfaces can realize the intermediate index value of two, suggesting a richer index spectrum when singularities are permitted.


\begin{remark}[Local computation along $\Gamma$]\label{rmk:local comp. on Gamma}
Fix a point on $\Gamma$ and choose local Fermi coordinates $(s,t)$ so that 
$\partial_s=\eta$ is the unit tangent to $\Gamma$ and $\partial_t=\tau_i$ is the inward conormal to the face $\Sigma_i$. 
At $t=0$, the coordinate vector fields commute: $[\partial_s,\partial_t]=0$. 
Set $D_i=\partial_t$ and write derivatives along $\Gamma$ as $('\!)=\partial_s$.

Differentiating $D_i(n_i^2)$ with respect to $s$ gives
\[
\partial_s(D_i(n_i^2)) = 2n_i' D_i n_i + 2n_i D_i n_i'.
\]
Therefore,
\[
n_i' D_i n_i + n_i D_i n_i' = \tfrac{1}{2}\,\partial_s(D_i(n_i^2)).
\]
Using the Frenet--Weingarten relations along $\Gamma$, 
$\nabla_\eta \eta = \kappa_i \tau_i$, 
we have
\[
n_i D_i n_i - \kappa_i n_i^2
= n_i' D_i n_i + n_i D_i n_i' - \kappa_i n_i^2.
\]
Combining the two identities and summing over $i=1,2,3$ yields the key relation
\[
\sum_{i=1}^3 (n_i D_i n_i - \kappa_i n_i^2)
= \tfrac{1}{2}\,\partial_s\!\Big(\sum_{i=1}^3 D_i(n_i^2)\Big),
\]
valid pointwise along $\Gamma$. 
Integrating over $\Gamma$ gives
\[
\sum_{i=1}^3 \int_\Gamma (n_i D_i n_i - \kappa_i n_i^2) = 0,
\]
since the right-hand side is the integral of a total tangential derivative. 
This proves the cancellation of the boundary terms used in the second part of the proof of Theorem~\ref{thm:main2}.
\end{remark}

\begin{proof}[Proof of Theorem \ref{thm:main2-n}]
We extend the proof of Theorem~\ref{thm:main2} to higher dimensions, highlighting the modifications needed for the general case.  \\ 
Let $YC \subset \mathbb{B}^n$ denote the flat $Y$-cone, obtained by gluing three planar half-disks along a common diameter at $120^\circ$ angles, as defined in \ref{def:Ycone}. Each face $\Sigma_j$ is flat, so the Jacobi operator on $\Sigma_j$ is
\[
J_j = \Delta_{\Sigma_j}, \quad |A_{\Sigma_j}|^2 = 0.
\]
The $J$-Steklov eigenvalue problem on each face depends only on the intrinsic two-dimensional geometry and the boundary conditions:
\begin{itemize}
    \item The eigenfunctions $w_k(r,\theta) = a_k r^k \cos(k\theta) + b_k r^k \sin(k\theta)$ have eigenvalues $\sigma_k = k$, independent of $n$.
    \item The compatibility condition along the junction, $f_1 + f_2 + f_3 = 0$ on $\Gamma$, is dimension-independent.
\end{itemize}

In the scalar (codimension one) case, the negative eigenspace of the index form
\[
Q(f,f) = \sum_{j=1}^3 \int_{\partial' \Sigma_j} (\sigma - 1) (f|_{\partial' \Sigma_j})^2
\]
is exactly two-dimensional, spanned by constant functions on the faces satisfying the compatibility condition.  

For $YC \subset \mathbb{R}^n$, the normal bundle has rank $n-2$. A direct computation using the explicit Steklov eigenfunctions shows that the two negative directions per normal component satisfy the junction compatibility independently and are linearly independent across the $n-2$ components. Therefore, the total Morse index is
\[
\operatorname{index}(YC) = 2(n-2).
\]

Now suppose $\Sigma \subset \mathbb{B}^n$ is a free-boundary minimal $Y$-surface with Morse index $2(n-2)$. Let $E_1,\dots,E_n$ denote the standard basis of $\mathbb{R}^n$, and let $NE_k$ be the normal projection of $E_k$ onto $\Sigma$. As in the three-dimensional case:
\begin{itemize}
    \item On each face $\Sigma_j$, $n^k_j = \langle E_k, N_j \rangle$ satisfies $J_j n^k_j = 0$.
    \item The index form satisfies $Q(NE_k,NE_k) < 0$ unless $NE_k \equiv 0$.
\end{itemize}
Since the Morse index is $2(n-2)$, at most $2(n-2)$ of the $NE_k$ can be linearly independent negative directions. Therefore, there exist $(n-2)$ independent parallel vector fields $E$ such that $NE \equiv 0$ on $\Sigma$, i.e., $(n-2)$ independent directions are everywhere tangent to $\Sigma$.  

This implies that $\Sigma$ is contained in a three-dimensional affine subspace $V \subset \mathbb{R}^n$.\\
Within $V$, $\Sigma$ is a two-dimensional free-boundary $Y$-surface in $\mathbb{B}^3$ with Morse index two. By Theorem~\ref{thm:main2}, the only such surface is the flat $Y$-cone. Hence $\Sigma$ coincides with the flat $Y$-cone in $V \subset \mathbb{R}^n$.
\medskip


Now, we compute the nullity is $5(n-2)$. We work facewise and componentwise in the normal bundle.  Let $\Sigma_1,\Sigma_2,\Sigma_3$ denote the three flat half-disk faces of $YC$.  Fix an orthonormal trivialization of the normal bundle on each face; because each face is flat and embedded as a planar half-disk, the Jacobi operator acts componentwise as the scalar Laplace operator $\Delta$ on each scalar normal component.  Thus a global normal Jacobi field may be written as a triple
\[
\Phi=(\phi_1,\phi_2,\phi_3),\qquad \phi_j:\Sigma_j\to\mathbb{R}^{n-2},
\]
with the understanding that $\phi_j$ is vector-valued (with $n-2$ scalar components) and each scalar component is harmonic on the interior of the corresponding half-disk.
\medskip

On each face $\Sigma_j$ the Jacobi equation $J\phi_j=0$ is the Laplace equation $\Delta \phi_j=0$ (componentwise). The linearized free boundary condition on the circular arc $\partial'\Sigma_j$ is exactly the Steklov condition
\[
\partial_\nu \phi_j = \phi_j \quad\text{on }\partial'\Sigma_j,
\]
where $\nu$ denotes the outward conormal of $\Sigma_j$ along the circular arc. Note that, this is the same boundary condition used in the index computation: the boundary term in the index form contains $(\sigma-1)\int_{\partial'\Sigma_j}(\phi_j)^2$ when a Steklov separation with eigenvalue $\sigma$ is used.  

Hence any scalar component of a global Jacobi field corresponds to a Steklov eigenfunction on the half-disk with eigenvalue $\sigma=1$.  It is classical, and follows by separation of variables or by homogeneity, that the Steklov eigenfunctions on the unit disk (and hence on the half-disk, restricted suitably) which have eigenvalue $\sigma=k$ are exactly the restrictions of homogeneous harmonic polynomials of degree $k$.  In particular, the $\sigma=1$ modes are precisely the degree-one harmonic polynomials, i.e., linear functions.  No higher degree ($k\ge2$) contributes to $\sigma=1$, and the constant functions correspond to $\sigma=0$.

Therefore every scalar component of a global Jacobi field on a face $\Sigma_j$ is necessarily the restriction of a linear function on that planar face.  Equivalently, in local coordinates $(s,t)$ on the half-disk where $s$ is the coordinate along the diameter $\Gamma$ and $t$ is the transverse coordinate (so $\Gamma=\{t=0\}$ and the circular arc is $s^2+t^2=1$, $t\ge0$), any scalar $\sigma=1$ solution on $\Sigma_j$ has the form
\[
f_j(s,t) = A_j s + B_j t,
\]
for constants $A_j,B_j\in\mathbb{R}$, depending on the chosen scalar component.
\medskip

Global admissibility across the three faces requires the usual junction compatibility condition along the common diameter $\Gamma$, namely
\begin{equation}\label{eq:compat-vector}
\phi_1|_{\Gamma} + \phi_2|_{\Gamma} + \phi_3|_{\Gamma} = 0
\end{equation}
as an equality of vector-valued functions along $\Gamma$.  Restricting the linear form $f_j(s,t)=A_j s + B_j t$ to $\Gamma$ (where $t\equiv0$) shows that the restriction depends only on the $A_j$ coefficient:
\[
f_j|_{\Gamma}(s) = A_j s.
\]
Hence the compatibility condition \eqref{eq:compat-vector} for a single scalar normal component reduces to the scalar identity
\[
(A_1 + A_2 + A_3)\, s \;=\; 0 \qquad\forall s\in[-1,1],
\]
and therefore
\[
A_1 + A_2 + A_3 = 0.
\]
When $\phi_j$ is vector-valued with $n-2$ scalar components, in the chosen normal frame, the compatibility condition is the vector equation obtained by applying the above scalar constraint componentwise; equivalently it imposes exactly $n-2$ independent scalar linear constraints (one per normal component) on the $3\times (n-2)$ collections of $A_j$-coefficients.
\medskip


For each face $j$ and for each scalar normal component there are two real parameters $(A_j,B_j)$ describing the degree-one harmonic solution on that face.  Thus for the three faces and a single scalar normal component there are $3\cdot 2=6$ parameters in total.  The compatibility condition along $\Gamma$ imposes one linear scalar constraint $A_1+A_2+A_3=0$, reducing the free parameters to $6-1=5$ for that scalar component.

Since the $n-2$ scalar normal components decouple (the Laplacian and the boundary condition act componentwise in the chosen trivialization), the same count applies independently to each normal component.  Therefore the full vector--valued kernel has dimension
\[
\underbrace{5}_{\text{degrees of freedom per scalar normal comp.}}\times \underbrace{n-2}_{\text{number of normal components}} \;=\; 5(n-2).
\]
\medskip


Finally, we must rule out any other possible kernel elements than those coming from the degree-one modes counted above.  Let $\Phi$ be any admissible global Jacobi field.  As observed in step (1), each scalar component of $\Phi$ is harmonic on the interior of each face and satisfies the Steklov condition $\partial_\nu \phi = \phi$ on the circular arcs; consequently each scalar component is a Steklov eigenfunction with eigenvalue $\sigma=1$ on each face.  By the classical homogeneous harmonic polynomial characterization of Steklov eigenfunctions, the only harmonic polynomials on the disk with Steklov eigenvalue $\sigma=1$ are linear polynomials.  Hence every scalar component is necessarily of the affine linear form $A_j s + B_j t$ on each face, and thus is included in the parameter family already counted.  No constants ($\sigma=0$) nor higher degree harmonics ($\sigma\ge2$) can contribute to the kernel.  This excludes any additional kernel elements.

\medskip

Combining the counts above yields $\dim\ker J=5(n-2)$.  That gives the general formula
\[
\dim\ker J = 5(n-2).
\]
\end{proof}



\begin{thebibliography}{15}

\bibitem{choe2025free}
J. ~Choe, 
\emph{Free boundary minimal surfaces and the reflection principle},
arXiv preprint arXiv:2501.02700, (2025).

%\bibitem{colding2011course}
%T. H. ~Colding  and W. P. ~Minicozzi,
%\emph{A course in minimal surfaces}, volume \textbf{121}.
%American Mathematical Soc., (2011).

\bibitem{devyver2019index}
B. ~Devyver,
\emph{Index of the critical catenoid},
Geometriae Dedicata, \textbf{199}:355--371, (2019).
Publisher: Springer.

\bibitem{fernandez2023free}
I.~Fernández, L.~Hauswirth, and P.~Mira,
\emph{Free boundary minimal annuli immersed in the unit ball},
 Arch. Ration. Mech. Anal., \textbf{247}(108), (2023).

\bibitem{fraser2014compactness}
A. ~Fraser and M. M. ~Li, 
\emph{Compactness of the space of embedded minimal surfaces with free boundary in three-manifolds with nonnegative Ricci curvature and convex boundary},
Journal of Differential Geometry, \textbf{96}(2):183--200, (2014).
Publisher: Lehigh University.

\bibitem{fraser2015uniqueness}
A. ~Fraser and R. ~Schoen, 
\emph{Uniqueness Theorems for Free Boundary Minimal Disks in
Space Forms}, International Mathematics Research Notices, 2015(\textbf{17}):8268--8274, (2015).
Publisher: Oxford University Press.

\bibitem{fraser2016sharp}
A. ~Fraser, and R. ~Schoen, 
\emph{Sharp eigenvalue bounds and minimal surfaces in the ball},
Inventiones mathematicae, \textbf{203}(3):823--890, (2016).
Publisher: Springer.

\bibitem{gilbarg1977elliptic}
D. ~Gilbarg and N. S.~Trudinger, 
\emph{Elliptic partial differential equations of second order}, volume \textbf{224}, number 2.
Springer, (1977).

\bibitem{matinpour2024morse}
E. ~Matinpour, 
\emph{Morse index of Y-singular minimal surfaces},
arXiv preprint arXiv:2410.15534, (2024).

\bibitem{mese2006parameterized}
C. ~Mese, and S. ~Yamada, 
\emph{The parameterized Steiner problem and the singular Plateau problem via energy},
In Amer. Math. Soc, (2006).
To appear in Trans.

\bibitem{mese2011local}
C. ~Mese and S. ~Yamada, 
\emph{Local uniformization and free boundary regularity of minimal singular surfaces},
Journal of Geometric Analysis, \textbf{21}(3):743--766, (2011).
Publisher: Springer.

\bibitem{nitsche1985stationary}
J. C. C. ~Nitsche, 
\emph{Stationary partitioning of convex bodies},
Archive for rational mechanics and analysis, \textbf{89}:1--19, (1985).
Publisher: Springer.

\bibitem{souam1997stability}
R. ~Souam, 
\emph{On stability of stationary hypersurfaces for the partitioning problem for balls in space forms},
Mathematische Zeitschrift, \textbf{224}(2):195--208, (1997).
Publisher: Berlin, J. Springer.


\bibitem{tran2016index}
H. ~Tran, 
\emph{Index characterization for free boundary minimal surfaces}, 
Communications in Analysis and Geometry, \textbf{28} (1), 189--222, (2020).


\bibitem{wang2021uniformization}
G. ~Wang, 
\emph{Uniformization of surfaces with boundary and the application to the triple junction surfaces with negative Euler characteristic},
arXiv preprint arXiv:2110.12656, (2021).

\bibitem{wang2022curvature}
G. ~Wang, 
\emph{Curvature estimates for stable minimal surfaces with a common free boundary},
Calculus of Variations and Partial Differential Equations, \textbf{61}(4):138, (2022).
Publisher: Springer.

\end{thebibliography}


\end{document}








\section{Main discussion and Rigidity Result}

Suppose that each face $\Sigma_j$ is minimally immersed in the unit ball $\mathbb{B}^n\subset\mathbb{R}^n$ through a smooth immersion
\[
\varphi_j:\Sigma_j\to\mathbb{B}^n .
\]
\medskip

Let $\Sigma_1, \Sigma_2,$ and $\Sigma_3$ be compact, smooth, two–dimensional, two-sided surfaces with smooth boundaries 
$\partial \Sigma_1, \partial \Sigma_2,$ and $\partial \Sigma_3$ lying in the unit ball $\mathbb{B}^n \subset \mathbb{R}^n$. 
Their boundary curves are allowed to meet the boundary sphere $\mathbb{S}^{n-1}$. Let $\Gamma$ be a smooth one-dimensional curve. 
For each $j = 1, 2, 3$, decompose the boundary of $\Sigma_j$ as
\[
\partial \Sigma_j = \partial' \Sigma_j \cup \partial'' \Sigma_j,
\]
where $\partial' \Sigma_j$ and $\partial'' \Sigma_j$ are disjoint smooth arcs (possibly sharing endpoints). 
Assume there exist smooth diffeomorphisms
\[
\phi_j : \Gamma \to \partial'' \Sigma_j
\]
that parameterize the boundary portions $\partial'' \Sigma_j$. 
We then \emph{identify} the three boundary arcs $\partial'' \Sigma_j$ via these maps; that is, for each $x \in \Gamma$, the points 
$\phi_1(x)$, $\phi_2(x)$, and $\phi_3(x)$ are regarded as a single point. 
After this identification, the three surfaces $\Sigma_1, \Sigma_2, \Sigma_3$ meet smoothly along the common curve $\Gamma$, 
forming a \emph{triple junction surface} where the three faces intersect along $\Gamma$.

In particular, we denote the resulting surface by
\[
\Sigma = \Big(\bigcup_{j=1}^3 \Sigma_j; \Gamma \Big),
\]
and refer to it as a \emph{triple junction surface with boundary}. 
The boundary of $\Sigma$ is the union of the non-identified boundary parts,
\[
\partial \Sigma = \bigcup_{j=1}^3 \partial' \Sigma_j,
\]
while the curve $\Gamma$ represents the identified portions of the boundaries, that is, the common junction curve along which the three faces are glued together.
\medskip


Suppose that each face $\Sigma_j$ is minimally immersed in the unit ball $\mathbb{B}^n\subset\mathbb{R}^n$ through a smooth immersion
\[
\varphi_j:\Sigma_j\to\mathbb{B}^n .
\]

%\begin{remark}
%\label{rem:gamma-embedded}
%In Definition~\ref{def: triple junction surface}, we allow each face $\Sigma_j$ to be an immersed minimal surface, so that $\Sigma_j$ may in principle have self–intersections. However, the junction curve $\Gamma$ along which the three faces meet at equal angles is necessarily \emph{embedded}. 

%Indeed, every point of $\Gamma$ is, by definition, a $Y$-point where exactly three smooth sheets meet at $120^\circ$. If $\Gamma$ were to self–intersect in $\mathbb{B}^n$, then at a self–intersection point $p$ there would be at least four distinct tangent sheets of $\Sigma$ meeting along $p$, which no longer forms a $Y$-type singularity. Such a point would constitute a higher–order or non-$Y$ singularity, contradicting the assumption that $\Sigma$ is a $Y$-surface. 

%Therefore, the local structure of $\Sigma$ near $\Gamma$ implies that $\Gamma$ must be an embedded curve in $\mathbb{B}^n$ and on $\mathbb{S}^{n-1}$.
%\end{remark}









In general, a function $u\in W^{1,2}_{com}(\Sigma)$ is called an eigenfunction associated with the eigenvalue $\lambda$ if, for every test function $v\in W^{1,2}_{com}(\Sigma)$, we have 
$$
\int_{\Sigma} \nabla u \cdot \nabla  vd\mu=\lambda \int_{\Sigma}uv\ d\mu.
$$
Assuming that $u$ is $C^2$ up to the boundary, integration by parts yields
$$
-\int_{\Sigma}v\Delta u\ d\mu + \int_{\Gamma} \sum_{i=1}^3 v_i \frac{\partial u_i}{\partial \tau_i} ds=\lambda \int_{\Sigma} uv\ d\mu ,
$$
for all $v \in W^{1,2}_{com}(\Sigma)$. In the interior of $\Sigma$, this implies the standard eigenvalue equation
$- \Delta u= \lambda u$.
However, on the common boundary $\Gamma$, the test functions satisfy the compatibility condition $\sum_{i=1}^3v_i=0$. Consequently, for the boundary term to vanish identically—i.e.,
$$
\int_{\Gamma}\sum_{i=1}^3 v_i\frac{\partial u_i}{\partial \tau_i} ds =0
$$ 
for all admissible $v$, it is necessary and sufficient that
$$
\frac{\partial u_i}{\partial \tau_i} = \frac{\partial u_j}{\partial \tau_j}
$$
on $\Gamma$ for all $i,j$.\\
This condition can be interpreted as requiring $C^1$-continuity of $u$ across the junction of the
three sheets of $\Sigma$.








